% !TeX root = ../manual.tex
% ===========================================================================
%  Introducción
% ===========================================================================

\section*{Acerca de este manual}
\addcontentsline{toc}{section}{Acerca de este manual}

Este manual explica, paso a paso, cómo usar \textbf{ReservaSport}, el sistema
de reserva de canchas deportivas del club. Cubre las dos formas de usar la
aplicación:

\begin{itemize}
  \item Como \textbf{usuario final}: la persona que consulta disponibilidad,
        reserva una cancha y gestiona sus propias reservas.
  \item Como \textbf{administrador}: la persona del club que gestiona el
        catálogo de canchas, las reservas de todos los usuarios, las cuentas
        y los reportes de uso.
\end{itemize}

Todas las capturas de este manual corresponden a la aplicación en
funcionamiento, con datos de ejemplo.

\subsection*{Requisitos}
\addcontentsline{toc}{subsection}{Requisitos}

Para usar ReservaSport solo hace falta un navegador web moderno (Chrome,
Edge, Firefox) y conexión a la dirección donde esté publicado el club. No
se necesita instalar ningún programa adicional.

\subsection*{Cómo está organizado este manual}
\addcontentsline{toc}{subsection}{Cómo está organizado este manual}

\begin{itemize}
  \item \textbf{\ref{sec:acceso} — Acceso al sistema}: crear una cuenta,
        iniciar sesión, el perfil y el cierre de sesión.
  \item \textbf{\ref{sec:usuario-final} — Guía para el usuario final}:
        consultar disponibilidad, reservar y cancelar.
  \item \textbf{\ref{sec:administrador} — Guía para el administrador}:
        el panel de administración completo.
  \item \textbf{\ref{sec:faq} — Preguntas frecuentes}: dudas comunes y qué
        hacer ante los mensajes más habituales del sistema.
\end{itemize}

\nota{A lo largo del manual, los nombres de botones y campos de la
aplicación aparecen en \boton{negrita} y los datos que escribes en
\campo{cursiva}, para que sea fácil ubicarlos en las capturas.}
